\section{Conclusion}
\label{sec:conclusion}

We have presented a four-chapter empirical diagnosis of EEG
foundation models under subject-level evaluation. Frozen EEG
FM representations are subject atlases, with subject identity
outweighing task-label variance by $10\times$ to $50\times$ in every
model$\times$dataset cell we examined. On the Komarov--Ko--Jung UCSD
Stress dataset, recalibrating the reported literature headline
($0.90$ BA) under subject-level CV yields a $0.52$--$0.58$ band, and a
controlled paired within-subject comparison against EEGMAT
rest-vs-arithmetic isolates the governing variable as the presence
or absence of a within-subject neural contrast anchor in the
published neuroscience literature. Seven architectures spanning five
orders of magnitude in parameter count share the same ceiling on the
contrast-bounded cohort; none escape it. Inside the ceiling,
pretrained FMs provide a modest but consistent $\sim\!5$\,pp
representation bump above classical baselines, but not a rescue. In
the anchored-contrast regime, the same FMs reach genuine rescue.

The Subject-Dominance Limit (SDL) diagnostic we derive from these
observations explains why several concurrent EEG-FM benchmark audits
report aggregate FM underperformance on cross-subject clinical
tasks, without requiring new architectures, losses, or datasets. The
operational deliverable is a pre-benchmark contrast-anchoring
protocol: practitioners can predict, from published neuroscience
literature alone, whether FM pretraining will rescue a target
clinical task, before committing to benchmark infrastructure.
Future pretraining objectives explicitly designed to reduce
subject-dominance in frozen representations are a falsifiable
direction that SDL predicts should produce the largest gains
specifically on contrast-bounded cohorts.
